\clearpage
\phantomsection% 
\addcontentsline{toc}{chapter}{ABSTRAK}
\thispagestyle{fancy}

\begin{center}
	\large \bfseries \MakeUppercase{Abstrak}\\
	\normalsize \normalfont {\thetitle}\\
	\normalsize \normalfont {\theauthor}\\
	\bigskip
	
	\normalsize \normalfont \justifying \singlespacing
	Kejadian jatuh pada lansia merupakan masalah serius yang memerlukan perhatian khusus dalam upaya pencegahan cedera. Menurut WHO, sekitar 37,3 juta kejadian jatuh memerlukan perhatian medis dan sekitar 684.000 kematian di seluruh dunia diakibatkan oleh jatuh. Ditambah lagi, Lansia di panti jompo memiliki risiko jatuh lebih tinggi dibanding lansia yang tinggal di rumah, dengan kejadian 30--50\% per tahun dan sekitar 40\% mengalami jatuh berulang. Permasalahan ini membutuhkan sistem berbasis model deteksi posisi jatuh yang cepat dan akurat, serta dapat bekerja 24 jam, terutama pada area yang tidak selalu diawasi untuk meningkatkan kualitas hidup lansia. Penelitian sebelumnya mengembangkan model deteksi jatuh berbasis visi dengan CNN dan LSTM, namun terdapat kelemahan pada kebutuhan daya komputasi yang tinggi untuk pengolahan citra. Penelitian ini bertujuan untuk mengembangkan Model deteksi berbasis \textit{Graph Convolutional Network} (GCN) menggunakan data \textit{pose landmark} yang secara teknik lebih ringan dari citra. Pada penelitian ini, telah dilakukan pengembangan model GCN dengan dua konfigurasi \textit{hyperparameter}, konfigurasi default dan konfigurasi hasil studi ablasi. Hasil eksperimen menunjukkan bahwa konfigurasi studi ablasi menghasilkan peningkatan kinerja yang signifikan dibandingkan dengan konfigurasi default. Secara keseluruhan (rerata), \textit{accuracy} meningkat sebanyak 0,018, \textit{precision} sebanyak 0,024, \textit{F1-score} sebesar 0,049 dan peningkatan terbesar terlihat pada \textit{recall}, sebesar 0,093. Hasil ini mengindikasikan bahwa konfigurasi studi ablasi dapat meningkatkan kinerja model GCN dalam mendeteksi kejadian jatuh pada lansia, yang memiliki potensi aplikasi yang besar dalam pengembangan sistem pemantauan kesehatan lansia yang berfokus pada jatuh dan tidak jatuh di masa depan.
    \vspace{20pt}
    
	\textbf{Kata Kunci: Lansia, Jatuh, Deep Learning, Graph Convolution Network, Mediapipe, Pose Landmark, studi ablasi}
	
	\vfill
	
\end{center}
\clearpage